
% ════════════════════════════════════════
% 模型的评价（07.模型评价.tex）
% ════════════════════════════════════════
\section{模型的评价}

\subsection{模型的优点}
\begin{itemize}[itemindent=2em]
\item \textbf{优点1：}线性规划模型结构简洁、约束刻画完整，所有决策变量均为连续变量，求解器可严格保证全局最优，避免了启发式算法陷入局部最优的风险。
\item \textbf{优点2：}滚动调整模型利用上境变量将分段惩罚费用精确线性化，在不引入整数变量的前提下完整刻画违约与超额的经济代价，计算效率高。
\item \textbf{优点3：}逐日滚动与储电量传递机制将全年大规模优化分解为数百个日尺度小问题，既降低了计算复杂度，又保持了全年储电量的连续性。
\item \textbf{优点4：}模型具有高度可迁移性，只需替换电价、负载与光伏序列即可适应不同季节、不同地区乃至波动电价场景，并通过灵敏度分析验证了稳健性。
\end{itemize}

\subsection{模型的缺点}
\begin{itemize}[itemindent=2em]
\item \textbf{缺点1：}模型假设各 10 分钟间隔内功率恒定，忽略了更短时间尺度上的功率波动对储能实际运行的影响。
\item \textbf{缺点2：}滚动调整模型假定负荷完全可预知、储能严格按计划执行，未考虑储能动态响应光伏实时偏差的闭环控制过程。
\end{itemize}
